\chapter{Conclusion}
\label{ch:conclusion}

This thesis asked how a researcher-operated adaptive reading platform can be architected so that sensing, eye-movement analysis, decision strategy, and intervention execution are independently replaceable, while supporting real Tobii-backed reading sessions and applying context-preserving micro-interventions that do not disrupt the participant's reading flow.
The motivation was concrete: prior Reading the Reader prototypes proved that gaze-driven adaptive reading works, but each was a monolithic fork that made the next study expensive to run, and no existing recorder or experiment builder combined real-time typographic intervention with pluggable, researcher-supervised decision logic.

In answer, we designed, built, and evaluated a two-screen adaptive reading platform.
A participant screen presents reflowable text that adapts to gaze, and a researcher console mirrors the participant's live gaze and gives the researcher control of the running session.
The adaptive loop is divided into four modules (sensing, analysis, decision, and intervention) behind stable contracts; interventions are applied manually or proposed by a pluggable decision provider in an advisory or autonomous mode; and every gaze sample, decision, and intervention is recorded into one authoritative session that can be exported, replayed, and re-imported.

The evaluation answered the four research questions with evidence rather than assertion.
Modularity (RQ1) holds where it is hardest to fake: the core's dependency rule is enforced by the build and a test, and three external providers connected without touching it.
The sensing pipeline (RQ2) ran at the tracker's rated \SI{90}{\hertz} with validity above 92\% and transport latency within the \SI{100}{\milli\second} budget for every sample.
Context preservation (RQ3) lowered both the reading-resume time and the post-intervention regression rate while recording an audited residual for every change.
Researcher control (RQ4) was sufficient to operate every session through one gated console and to reconstruct all reported results from the exported records alone.
Two findings we state without softening: the automated decision path was exercised within budget in a single validation run rather than at campaign scale, and the context-preserving restore moved the reader too far on small reflows, a flaw we revised to hold the captured offset and have so far confirmed only geometrically.

The lasting contribution is the design knowledge the platform embodies, distilled into four principles that outlive this implementation: enforce a module boundary as a build-time invariant, validate an integration contract from outside rather than only by its authors, measure the cost of each quality attribute the architecture claims instead of asserting it, and separate a change's decision from its commit so a disruptive update lands at a natural break.
With reproducible, replayable records now in hand and a decision boundary open to an external provider, the platform supplies the two things the project's envisaged classifier has always presupposed but never had: data of a quality it can be trained on, and a live loop in which it can be connected, run, and judged.
The controlled efficacy study and the campaign-scale, AI-driven decision-making this thesis deliberately left out of scope thus become configuration choices for the work that follows, not rewrites of core code.
